\chapter{Lumpectomy}

\section*{What Is This Surgery?}
A lumpectomy, also known as breast-conserving surgery (BCS) or wide local excision, is a surgical procedure designed to remove a malignant tumor from the breast while preserving the majority of the breast tissue and its overall shape. This procedure involves the precise excision of the breast cancer along with a surrounding margin of healthy, normal tissue to ensure complete removal of all macroscopic and microscopic cancerous cells. It is a cornerstone of treatment for early-stage breast cancer, typically followed by adjuvant radiation therapy to the remaining breast tissue to significantly reduce the risk of local recurrence. The primary goal is to achieve oncologic control equivalent to mastectomy while offering superior cosmetic and psychological outcomes for suitable patients, thereby maintaining quality of life. This approach is a critical component of a multidisciplinary treatment plan for breast cancer, integrating surgical, radiation, and medical oncology expertise.

\section*{Alternative Names}
\begin{itemize}
  \item Breast-Conserving Surgery (BCS)
  \item Wide Local Excision (WLE)
  \item Partial Mastectomy
  \item Segmental Mastectomy
  \item Quadrantectomy (when a full quadrant of breast tissue is removed)
  \item Tylectomy (a less common term for tumor excision)
\end{itemize}

\section*{Relevant Surgical Anatomy}
The breast is a modified sweat gland situated on the anterior chest wall, extending from the second to the sixth rib vertically and from the sternal edge to the mid-axillary line horizontally. It is composed of 15-20 lobes of glandular tissue, each draining into a lactiferous duct that opens onto the nipple. These lobes are embedded within a fibrous and adipose stroma, supported by suspensory ligaments known as Cooper's ligaments, which attach to the overlying skin and underlying pectoral fascia. The breast is divided into four quadrants (upper outer, upper inner, lower outer, lower inner) and a central subareolar region, with the upper outer quadrant extending into the axilla as the axillary tail of Spence, a common site for breast cancer.

The arterial supply to the breast is rich, primarily from perforating branches of the internal mammary artery (medially), the lateral thoracic artery and thoracoacromial artery (laterally), and intercostal arteries (posteriorly). Venous drainage largely parallels the arterial supply, emptying into the axillary vein, internal mammary vein, and intercostal veins. Lymphatic drainage is of paramount surgical importance; approximately 75\% of lymph from the breast drains to the axillary lymph nodes (levels I, II, III), with the remaining drainage to the internal mammary nodes, supraclavicular nodes, and contralateral breast. Key nerves in the vicinity include the intercostobrachial nerve (sensory to the medial upper arm), and the long thoracic and thoracodorsal nerves (motor to serratus anterior and latissimus dorsi, respectively), which are at risk during axillary dissection. Surgical landmarks include the nipple-areola complex, inframammary fold, and the palpable border of the pectoralis major muscle.

\section{Surgical Principle \& Rationale}
The fundamental principle of lumpectomy is the complete surgical removal of the primary breast tumor with a surrounding rim of uninvolved, healthy breast tissue, known as a "clear margin," while preserving as much normal breast parenchyma as possible. This approach aims to eradicate all macroscopic and microscopic cancer cells from the surgical site, thereby minimizing the risk of local recurrence within the breast. The excised specimen is meticulously oriented by the surgeon using sutures or clips to allow the pathologist to accurately assess the three-dimensional margins and determine if any cancer cells are present at the edges of the removed tissue. Achieving clear margins is paramount for oncologic safety.

Concurrently, a sentinel lymph node biopsy (SLNB) is often performed to assess for regional lymph node metastasis. This involves injecting a radioactive tracer and/or blue dye near the tumor, which then travels to the first few lymph nodes (sentinel nodes) that drain the tumor area. These nodes are surgically removed and examined for cancer cells. If the sentinel nodes are negative, further axillary dissection is typically avoided, reducing the risk of lymphedema and other axillary morbidity. The overarching technical goals are to achieve clear surgical margins, accurately stage the axilla, and reconstruct the breast defect using oncoplastic techniques to maintain an acceptable cosmetic outcome, often utilizing local tissue rearrangement or reduction mammoplasty principles to reshape the remaining breast tissue and ensure symmetry.

\section{History \& Surgical Evolution}
The concept of breast-conserving surgery represents a significant evolution in breast cancer treatment, moving away from the radical mastectomy that dominated the 20th century, championed by William Halsted. Early pioneers like Dr. George Crile Jr. in the 1960s challenged the prevailing radical approach, advocating for less extensive surgery for selected patients based on the understanding of breast cancer as a systemic disease rather than purely local. However, widespread acceptance came with the publication of landmark randomized controlled trials in the 1980s.

The National Surgical Adjuvant Breast and Bowel Project (NSABP) B-06 trial, led by Dr. Bernard Fisher and published in 1985, demonstrated unequivocally that lumpectomy followed by radiation therapy offered equivalent long-term survival rates to total mastectomy for women with early-stage breast cancer. Similar findings emerged from the Milan I trial, conducted by Dr. Umberto Veronesi, which also showed comparable long-term survival for quadrantectomy with axillary dissection and radiation versus radical mastectomy. These pivotal studies revolutionized breast cancer management, establishing breast-conserving surgery as a safe and effective alternative to mastectomy for appropriately selected patients. Subsequent advancements have focused on refining margin definitions, improving localization techniques for non-palpable lesions (e.g., wire, radioactive seed, magnetic seed), and integrating oncoplastic surgery to enhance cosmetic outcomes while maintaining oncologic safety.

\section{Clinical Indications}
\begin{itemize}
  \item To definitively treat early-stage invasive breast cancer, typically Stage I or II, where the tumor size is amenable to complete removal with clear margins while preserving breast cosmesis.
  \item For ductal carcinoma in situ (DCIS), which is a non-invasive form of breast cancer, to remove the abnormal cells and prevent progression to invasive disease, particularly for unifocal lesions.
  \item When the patient expresses a strong desire for breast preservation and meets the oncologic criteria for breast-conserving therapy, including the ability to undergo adjuvant radiation.
  \item As part of a neoadjuvant treatment strategy, where chemotherapy or hormone therapy is given before surgery to shrink a larger tumor, making it suitable for lumpectomy that would otherwise require mastectomy.
  \item For small, localized recurrences of breast cancer after prior breast-conserving surgery, provided the initial radiation field is not compromised and clear margins can be achieved.
  \item When the tumor-to-breast size ratio allows for adequate tumor excision with clear margins and an acceptable cosmetic result, often guided by oncoplastic principles.
  \item In cases where multidisciplinary tumor board review recommends breast-conserving surgery based on favorable tumor biology, patient factors, and the availability of adjuvant radiation therapy.
  \item For certain benign but high-risk lesions, such as atypical ductal hyperplasia (ADH) or lobular carcinoma in situ (LCIS), when an excisional biopsy is performed to rule out occult malignancy or for definitive removal.
\end{itemize}

\section*{Clinical Positioning}
Lumpectomy is predominantly considered a primary surgical treatment for early-stage breast cancer. For many women diagnosed with ductal carcinoma in situ (DCIS) or invasive breast cancer stages I and II, it serves as the initial and definitive surgical intervention, often followed by adjuvant radiation therapy and systemic treatments. Its role as a primary procedure is contingent upon the tumor's characteristics, such as size, location, multifocality, and biological subtype, as well as patient factors like breast size, overall health, and personal preference for breast preservation.

In certain scenarios, lumpectomy may also function as a secondary or salvage procedure. This occurs most commonly when initial surgical margins are found to be positive after a previous lumpectomy, necessitating a re-excision to achieve clear margins. This re-excision is still a breast-conserving approach. Less frequently, it might be considered for a local recurrence after a prior lumpectomy, though mastectomy is often preferred in such cases, especially if prior radiation therapy limits further breast conservation. In the context of neoadjuvant therapy, lumpectomy becomes the primary surgical procedure after the tumor has been downstaged, allowing for breast preservation where it might not have been possible initially.

\section{Surgical Technique \& Procedure}
The lumpectomy procedure is performed under general anesthesia, with the patient typically positioned supine with the ipsilateral arm abducted on an arm board to allow access to the axilla. Preoperatively, for non-palpable lesions, the tumor's exact location is often marked by a radiologist using imaging guidance (e.g., mammography, ultrasound, or MRI) to place a localization wire, radioactive seed, or magnetic seed directly into or adjacent to the lesion. This ensures precise intraoperative identification.

The surgeon then makes an incision, ideally placed in a cosmetically favorable location such as along the areolar border, in the inframammary fold, or following Langer's lines, to minimize visible scarring. Through this incision, the tumor, along with the localization device and a surrounding margin of healthy breast tissue, is carefully excised. The specimen is meticulously oriented by the surgeon using sutures or clips to indicate specific margins (e.g., superior, inferior, medial, lateral, deep, superficial) for the pathologist. Intraoperative specimen radiography or ultrasound may be used to confirm removal of the lesion and assess initial margin adequacy.

Following tumor excision, if indicated, a sentinel lymph node biopsy is performed through a separate small incision in the axilla to identify and remove the first draining lymph nodes for pathological assessment. Hemostasis is achieved using electrocautery. The breast cavity may then be reshaped using oncoplastic techniques, such as local tissue rearrangement (e.g., glandular flaps, rotation flaps) or reduction mammoplasty principles, to minimize deformity and ensure breast symmetry. Finally, the incision is closed in layers, often with absorbable sutures for the deeper tissues and subcuticular sutures for the skin, and a sterile dressing is applied. Drains are typically not used unless a significant dead space or extensive axillary dissection was performed.

\section{Preoperative Workup \& Preparation}
Thorough preoperative preparation is crucial for a safe and successful lumpectomy, ensuring both oncologic efficacy and patient safety. This typically begins with a comprehensive diagnostic workup, including mammography, ultrasound, and often breast MRI, to fully characterize the tumor, assess its extent, and rule out multifocal or multicentric disease. A biopsy (core needle or excisional) must confirm the diagnosis of malignancy and provide tumor biology markers (ER, PR, HER2). Patients will undergo a complete medical history and physical examination, along with routine preoperative laboratory tests such as a complete blood count, coagulation panel, and basic metabolic panel to assess overall health and fitness for surgery.

\begin{itemize}
  \item \textbf{Anesthesia Consultation:} Evaluation by an anesthesiologist to assess fitness for general anesthesia, review current medications, and discuss pain management strategies for the perioperative period.
  \item \textbf{Medication Review:} Patients are instructed to discontinue blood-thinning medications (e.g., aspirin, NSAIDs, warfarin, novel oral anticoagulants) for a specified period before surgery, as advised by the surgeon and anesthesiologist, to minimize bleeding risk.
  \item \textbf{NPO Status:} Patients must adhere to "nil per os" (NPO) guidelines, meaning no food or drink for a specified number of hours (typically 6-8 hours) before surgery to prevent aspiration during anesthesia induction.
  \item \textbf{Preoperative Imaging/Localization:} If the tumor is non-palpable, a radiologist will perform a wire localization, radioactive seed placement, or magnetic seed placement on the day of or day before surgery to guide the surgeon.
  \item \textbf{Informed Consent:} Detailed discussion with the surgeon regarding the procedure, potential risks, benefits, alternatives (e.g., mastectomy), expected outcomes, and the possibility of re-excision or conversion to mastectomy, followed by signing of consent forms.
  \item \textbf{Antibiotic Prophylaxis:} Prophylactic antibiotics (e.g., cefazolin) are typically administered intravenously within 60 minutes prior to incision to reduce the risk of surgical site infection.
  \item \textbf{Discussion of Adjuvant Therapy:} Patients should have a preliminary understanding of potential post-operative treatments, such as radiation therapy, chemotherapy, or hormone therapy, which will be determined by the final pathology results and multidisciplinary tumor board recommendations.
\end{itemize}

\section*{Contraindications \& Precautions}
\textbf{Absolute Contraindications:}
\begin{itemize}
  \item \textbf{Multicentric Breast Cancer:} Presence of two or more distinct tumor foci in different quadrants of the breast, making complete excision with a single lumpectomy impossible without significant cosmetic compromise or high recurrence risk.
  \item \textbf{Diffuse Malignant Microcalcifications:} Widespread calcifications indicative of extensive ductal carcinoma in situ (DCIS) that cannot be adequately removed by lumpectomy with clear margins.
  \item \textbf{Persistent Positive Margins:} Inability to achieve clear surgical margins after multiple attempts at re-excision, indicating diffuse disease not amenable to breast conservation.
  \item \textbf{Prior Therapeutic Radiation to the Breast/Chest Wall:} Previous radiation therapy to the same breast significantly increases the risk of severe radiation reactions, poor wound healing, and recurrence with further radiation.
  \item \textbf{Inflammatory Breast Cancer:} This aggressive form of cancer typically requires mastectomy due to its diffuse nature, dermal lymphatic invasion, and high risk of local recurrence.
  \item \textbf{Large Tumor in a Small Breast:} When the tumor-to-breast size ratio is so unfavorable that lumpectomy would result in severe cosmetic deformity or an unacceptably small remaining breast, even with oncoplastic techniques.
\end{itemize}

\textbf{Relative Contraindications and Precautions:}
\begin{itemize}
  \item \textbf{Connective Tissue Diseases:} Conditions like scleroderma, systemic lupus erythematosus, or other autoimmune disorders can predispose patients to severe radiation reactions, poor wound healing, and increased fibrosis, making adjuvant radiation therapy challenging.
  \item \textbf{Pregnancy:} While lumpectomy can be performed in the second or third trimester, radiation therapy is contraindicated during pregnancy, delaying a critical component of breast conservation and potentially impacting fetal development.
  \item \textbf{Very Large Tumor Size:} Although neoadjuvant therapy can shrink tumors, very large tumors may still result in poor cosmetic outcomes even after lumpectomy, even if clear margins are achieved.
  \item \textbf{Patient Preference:} Some patients may simply prefer mastectomy for peace of mind, to avoid radiation therapy, or due to anxiety about recurrence, even if oncologically eligible for lumpectomy.
  \item \textbf{Unfavorable Tumor Biology:} Certain aggressive tumor types (e.g., triple-negative breast cancer) might prompt a more cautious discussion about the benefits of breast conservation versus mastectomy, though this is less common as an absolute contraindication to lumpectomy itself.
  \item \textbf{Lack of Access to Radiation Therapy:} If a patient cannot access or tolerate adjuvant radiation therapy due to geographical, financial, or medical reasons, lumpectomy alone may not be an appropriate oncologic treatment, as it significantly increases local recurrence risk.
\end{itemize}

\section*{Complications \& Risks}
As with any surgical procedure, lumpectomy carries inherent risks, which can be broadly categorized as general surgical risks and procedure-specific complications.

\begin{itemize}
  \item \textbf{General Surgical Risks:}
    \begin{itemize}
      \item \textbf{Anesthesia Complications:} Reactions to anesthetic agents, nausea, vomiting, respiratory issues, or, rarely, more severe cardiac or neurological events such as stroke or myocardial infarction.
      \item \textbf{Bleeding:} Intraoperative or postoperative hemorrhage, which may necessitate transfusion, re-operation for hematoma evacuation, or prolonged recovery.
      \item \textbf{Infection:} Surgical site infection (SSI), characterized by redness, swelling, pain, warmth, or purulent drainage, potentially requiring antibiotics, wound drainage, or further intervention.
      \item \textbf{Deep Vein Thrombosis (DVT) / Pulmonary Embolism (PE):} Formation of blood clots in the legs that can travel to the lungs, a rare but serious complication, mitigated by early ambulation and, in some cases, prophylactic anticoagulation.
    \end{itemize}
  \item \textbf{Procedure-Specific Risks:}
    \begin{itemize}
      \item \textbf{Seroma:} Accumulation of clear, yellowish fluid in the surgical cavity, which may cause swelling and discomfort and often requires aspiration in the outpatient setting.
      \item \textbf{Hematoma:} Collection of blood in the surgical site, potentially causing significant pain, swelling, bruising, and skin discoloration, and sometimes requiring surgical drainage.
      \item \textbf{Cosmetic Deformity:} Changes in breast shape, size, or symmetry, including indentation, puckering, nipple distortion, or breast asymmetry, which can range from subtle to significant and may require revision surgery.
      \item \textbf{Positive Margins:} Microscopic cancer cells found at the edge of the excised tissue, necessitating a re-excision (second surgery) or, in some cases, conversion to mastectomy to achieve clear margins.
      \item \textbf{Persistent Pain:} Chronic pain or discomfort at the surgical site or in the axilla, sometimes related to nerve irritation or neuroma formation, which can be debilitating.
      \item \textbf{Nerve Injury:} Damage to sensory nerves, particularly the intercostobrachial nerve in the axilla during sentinel lymph node biopsy, leading to numbness, tingling, or pain in the upper arm and axilla.
      \item \textbf{Lymphedema:} Swelling of the arm, hand, or breast due to impaired lymphatic drainage, more common and severe after axillary lymph node dissection but can occur rarely after sentinel lymph node biopsy.
      \item \textbf{Fat Necrosis:} Hardening or lumpiness in the breast tissue due to the death of fat cells, which can sometimes be mistaken for recurrence on imaging or physical exam.
      \item \textbf{Delayed Wound Healing:} Slower than expected healing of the incision, potentially leading to wound dehiscence, prolonged recovery, or increased risk of infection.
      \item \textbf{Scarring:} Formation of a visible scar, which may be hypertrophic or keloid in susceptible individuals, impacting cosmetic outcome.
      \item \textbf{Skin Discoloration:} Bruising and temporary or permanent changes in skin pigmentation around the surgical site.
    \end{itemize}
\end{itemize}

\section*{Postoperative Recovery Timeline}
Immediately after a lumpectomy, patients are transferred to the Post-Anesthesia Care Unit (PACU) for close monitoring of vital signs, pain levels, and recovery from general anesthesia. Pain management typically involves oral analgesics, and anti-nausea medication may be administered as needed. Most patients are discharged home the same day or the following morning, depending on their recovery, the extent of surgery (e.g., if axillary dissection was performed), and any co-morbidities. Before discharge, patients receive detailed instructions on wound care, activity restrictions, pain medication regimen, and signs of potential complications to watch for.

During the first 1-2 weeks post-procedure, patients are advised to avoid strenuous activities, heavy lifting (typically >5-10 lbs), and excessive arm movements on the operative side to promote healing and minimize swelling and seroma formation. A small amount of bruising, swelling, and mild discomfort around the incision site is normal. If a drain was placed (uncommon for lumpectomy alone, more so with axillary dissection), patients receive instructions on drain care and output measurement, with removal typically occurring when output is minimal. The first follow-up appointment with the surgeon typically occurs within 1-2 weeks to check the incision, remove any non-absorbable sutures or clips, and discuss the preliminary pathology results. This appointment is crucial for planning subsequent adjuvant therapies, such as radiation, chemotherapy, or hormone therapy, which usually commence several weeks after surgery once the surgical site has sufficiently healed. Full recovery of normal activities and exercise typically takes 4-6 weeks.

\section*{Surgical Findings \& Pathology}
During a lumpectomy, the surgeon documents gross findings such as the location and size of the tumor, its relationship to the localization device, and the appearance of the surrounding breast tissue. Any suspicious lymph nodes identified during sentinel lymph node biopsy are also noted. The excised breast tissue specimen is meticulously oriented by the surgeon using sutures or clips to mark specific margins (e.g., superior, inferior, medial, lateral, deep, superficial) before being sent to pathology. This precise orientation is critical for the pathologist to accurately assess the three-dimensional margins.

The pathology report on the resected tissues provides definitive diagnostic and prognostic information. Key elements include the precise type of cancer (e.g., invasive ductal carcinoma, invasive lobular carcinoma, ductal carcinoma in situ), its histological grade (indicating aggressiveness), and the exact tumor size. Crucially, the report details the status of the surgical margins, indicating whether the tumor was completely removed with a clear rim of healthy tissue (negative margins, typically >1-2 mm) or if cancer cells are present at the edge of the specimen (positive margins). The report also assesses for lymphovascular invasion, which indicates cancer cells in blood or lymphatic vessels. For invasive cancers, immunohistochemical staining results for hormone receptors (estrogen receptor - ER, progesterone receptor - PR) and HER2 status are vital for determining eligibility for hormone therapy and targeted therapies. Finally, if a sentinel lymph node biopsy was performed, the report will indicate the number of lymph nodes examined and whether any contain metastatic cancer cells, which is a major prognostic factor and influences the need for further axillary surgery or radiation.

\section*{Expected Postoperative Course}
A normal, successful postoperative course following a lumpectomy involves a relatively smooth recovery with minimal complications. Patients can expect mild to moderate pain at the surgical site, which is typically well-controlled with oral analgesics and gradually resolves over several days to a week. Some bruising and swelling around the incision are common and will subside within 1-2 weeks. The incision should heal cleanly, with sutures or staples (if used) removed at the first follow-up visit. Patients should experience a gradual return to normal arm and shoulder function, with most able to resume light activities within a few days and more strenuous activities within 4-6 weeks. Cosmetically, the goal is to preserve the natural contour and appearance of the breast as much as possible, with any minor indentations or asymmetries often improving over several months. The absence of fever, excessive redness, purulent drainage, or significant new lumps indicates an uncomplicated healing process, setting the stage for subsequent adjuvant therapies.

\section{Postoperative Care \& Follow-up}
Following a lumpectomy, a structured follow-up plan is essential for monitoring recovery, managing adjuvant therapies, and detecting any potential recurrence.
\begin{itemize}
  \item \textbf{Postoperative Surgical Visits:} Typically 1-2 weeks after surgery for wound check, removal of non-absorbable sutures or clips, and detailed discussion of the final pathology report and its implications.
  \item \textbf{Oncology Consultations:} Referrals to a medical oncologist for systemic therapy planning (chemotherapy, hormone therapy, targeted therapy) and a radiation oncologist for adjuvant radiation therapy to the breast, based on pathology results and tumor characteristics.
  \item \textbf{Radiation Therapy:} Usually begins 3-6 weeks after surgery, once the surgical site has healed, and typically lasts for 3-6 weeks, depending on the regimen (e.g., whole breast radiation, partial breast radiation).
  \item \textbf{Imaging Surveillance:} Annual mammograms of both breasts are recommended, often starting 6-12 months after surgery, to monitor for recurrence in the treated breast or new cancers in either breast.
  \item \textbf{Physical Therapy/Rehabilitation:} May be recommended for patients experiencing shoulder stiffness, reduced range of motion, or mild lymphedema, especially if axillary dissection was performed, to restore full function.
  \item \textbf{Activity Resumption:} Gradual return to normal activities, with specific guidance on when to resume strenuous exercise or heavy lifting, usually over several weeks, tailored to individual recovery.
  \item \textbf{Warning Signs:} Patients are educated on warning signs to report immediately, such as fever, increasing redness or swelling, pus from the incision, new lumps in the breast or axilla, persistent severe pain, or unexplained arm swelling.
\end{itemize}
Long-term follow-up with the oncology team is crucial for ongoing surveillance, management of potential side effects from adjuvant therapies, and addressing any concerns about recurrence or new primary cancers.

\section*{Epidemiology}
Breast cancer is the most common cancer among women worldwide, excluding non-melanoma skin cancers, and a leading cause of cancer-related death. In the United States, approximately 1 in 8 women (13\%) will develop invasive breast cancer over their lifetime. Lumpectomy, as a primary surgical treatment, is performed in a significant proportion of these cases, particularly for early-stage disease. The incidence of breast cancer increases with age, with the median age at diagnosis being around 62 years, though it can occur at any age. While predominantly affecting women, breast cancer can also occur in men, though it is rare (less than 1\% of all breast cancers). The frequency of lumpectomy versus mastectomy has shifted over decades, with breast-conserving surgery becoming increasingly common due to robust evidence of equivalent oncologic outcomes and improved quality of life. Currently, lumpectomy with radiation therapy is the preferred surgical approach for about 60-70\% of women with early-stage breast cancer in many developed countries. Mortality rates for early-stage breast cancer treated with lumpectomy and adjuvant therapy are low, with 5-year survival rates exceeding 90\%. Morbidity associated with lumpectomy is generally low, primarily involving local complications such as seroma, hematoma, and cosmetic changes, with severe complications being rare.

\section*{Outcomes \& Prognosis}
The outcomes and prognosis following a lumpectomy for early-stage breast cancer are generally excellent, particularly when combined with adjuvant radiation therapy and appropriate systemic treatments. For patients with Stage I or II breast cancer, long-term survival rates are comparable to those achieved with mastectomy, as demonstrated by landmark trials like NSABP B-06 (Fisher et al.) and the Milan I trial (Veronesi et al.). Five-year survival rates often exceed 90\%, and 10-year survival rates remain high, reflecting effective local and systemic control. Local recurrence rates in the treated breast are typically low, ranging from 5-10\% over 10 years when lumpectomy is followed by radiation therapy; without radiation, local recurrence rates are significantly higher.

Prognostic factors that significantly influence outcomes include tumor size, histological grade, lymph node status (especially sentinel lymph node involvement), hormone receptor status (estrogen and progesterone receptor positivity), and HER2 status. Tumors that are small, low-grade, node-negative, hormone receptor-positive, and HER2-negative generally have the most favorable prognosis. Functional outcomes are typically very good, with most patients experiencing a full recovery of arm and shoulder function. Cosmetic outcomes vary but are often satisfactory, especially with modern oncoplastic techniques, contributing positively to quality of life. While lumpectomy aims for local control, the overall prognosis is also heavily influenced by the risk of distant metastasis, which is addressed through systemic therapies tailored to the tumor's biology and patient characteristics.

\section*{References}
\begin{enumerate}
  \item MedlinePlus. Lumpectomy. U.S. National Library of Medicine; 2024. Available from: \url{https://medlineplus.gov/lumpectomy.html}
  \item Bland KI, Copeland EM. The Breast: Comprehensive Management of Benign and Malignant Diseases. 5th ed. Elsevier; 2018.
  \item National Comprehensive Cancer Network (NCCN). NCCN Clinical Practice Guidelines in Oncology: Breast Cancer. Version 1.2024. Available from: \url{https://www.nccn.org/guidelines/guidelines-detail?category=1&id=1419}
  \item Fisher B, Redmond C, Poisson R, et al. Ten-year results of a randomized clinical trial comparing radical mastectomy and total mastectomy with or without radiation therapy with breast-conserving surgery and radiation therapy in the treatment of node-negative breast cancer. N Engl J Med. 1985;312(11):665-673.
\end{enumerate}

\clearpage